\section{Introduction}
Legged platforms are increasingly deployed for inspection and search tasks in environments
that defeat wheeled locomotion: collapsed structures, forested slopes, and industrial sites
mid-renovation. In these settings the dominant failure mode is not gross planning error but
local foothold mis-placement -- a leg lands on a loose rock, an unsupported tile edge, or a
root hidden beneath leaf litter, and the resulting slip propagates into a stumble or fall.
Elevation-map-based foothold planners~\cite{singh2018survey,tanaka2023perception} handle
open, well-observed terrain well, but degrade sharply near occlusion boundaries, where a
single missing depth return can shift the estimated ground height by several centimeters.

Prior work addresses this gap along two lines. The first improves the map itself, fusing
multiple viewpoints or temporal windows to fill occluded regions~\cite{tanaka2023perception}.
The second bypasses maps for reactive, force-driven recovery after a stumble has already
begun~\cite{oduya2020reactive}. Both are valuable but incomplete: better maps still contain
residual noise near clutter, and reactive recovery only engages after contact quality has
already degraded. We target the interval between the two -- refining a foothold using recent
contact history \emph{before} the leg commits to stance.

This paper contributes: (i) \textsc{ContactRefine}, a foothold refinement module that jointly
conditions on local elevation and a short window of per-leg contact force history; (ii) a
contrastive training objective that requires no manual terrain labels, using slip and
force-spike events recorded during teleoperated data collection as implicit supervision; and
(iii) a $1{,}400$-trial physical evaluation across five terrain classes showing consistent
stumble-rate and cost-of-transport improvements with negligible added latency.
